% ===================================================================
%  اللوحة نمرة ٤ — سنّ النضج والشيخوخة.
%
%  Traduction, faite en 2026, en arabe egyptien ecrit, sur l'ido de
%  texto/io. Regles de la colonne : voir 10-lawha-01.tex. Deux
%  scenes, deux numerotations : les cles portent c1 et c2.
%
%  LE SUFFIXE DE METIER -جي, VENU DU TURC, EST A L'EGYPTIEN CE QUE 佬
%  EST AU CANTONAIS. Il apparait ici pour la premiere fois :
%
%    صفرجي (11)   le domestique de table  — de « sofra », la table
%
%  et il court ensuite dans toute la langue : بوسطجي le facteur,
%  عربجي le cocher, قهوجي le cafetier, مكوجي le repasseur, جزمجي le
%  cordonnier, كبابجي le rotisseur. L'arabe standard n'a pas ce
%  procede : il faut chaque fois un participe ou une annexion. C'est
%  exactement le rapport qu'on avait releve entre 泥水佬 et 泥水匠.
%
%  LA MEDECINE ET LES PAPIERS SONT EGYPTIENS ET NE SE DISENT PAS
%  AUTREMENT :
%
%    أجزخانة (16)  la pharmacie   MSA : صيدلية  — du turc eczahane
%    روشتّة (27)    l'ordonnance   MSA : وصفة   — de l'italien ricetta
%    نتيجة (43)     le calendrier  MSA : تقويم
%    نضّارة         les lunettes
%    سخونية        la fievre      MSA : حمّى
%    كحّة           la toux        MSA : سعال
%    شربة (28)     la potion
%    شغّالة (14)    la servante    MSA : خادمة
%
%  ET LE MOBILIER DE LA MAISON BOURGEOISE PARLE FRANCAIS ET ITALIEN :
%  فوتيه (34) le fauteuil, فيرندة (8) la veranda, بالطو (7) le
%  manteau, جاكيتة (11) la jaquette, منشون (13) le manchon, روب (35)
%  la robe de chambre, أزازة (52) la bouteille, بريمة (51) le
%  tire-bouchon, فلّينة (53) le bouchon, تورتة (46) le gateau.
%
%  LA FAMILLE, ELLE, EST TOUTE EGYPTIENNE : فرح la noce, معازيم (22)
%  les invites, قرايب les parents, نسيب (30) le beau-frere, تيتة (29,
%  23) la grand-mere, جدو (33) le grand-pere. Ces deux derniers sont
%  les mots que les enfants emploient, et l'ido dit « avino » et
%  « avulo » sans plus de ceremonie.
%
%  ET LA NEIGE EST تلج (19), non ثلج : la meme regle qu'aux tableaux
%  2 et 3, appliquee cette fois a l'hiver. De meme ديل la queue, non
%  ذيل, et ضلمة l'obscurite, non ظلمة.
% ===================================================================

%%K t04-apar-1 apar
\VUsaut{4.0mm}
\VUtitre{15.0pt}{60}{اللوحة نمرة ٤}
\VUsaut{1.6mm}
\VUtitre{11.4pt}{0}{سنّ النضج والشيخوخة.}
\VUsaut{3.2mm}

%%K t04-tit-1 sub
\VUcentre{11.4pt}{0}{\textit{المشهد الأول.}}
\VUsaut{1.6mm}
\VUtitre{11.4pt}{0}{عزومة الفرح.}
\VUsaut{2.6mm}

%%K t04-tit-2 sub
\VUpk{10.2pt}{40}{الخريف.}
\VUsaut{3.0mm}

%%K t04-c1-01-1 p
1. --- الشباب كبروا. وواحد فيهم، يوحنا، بيتجوّز. وإحنا حاضرين
\VUgras{عزومة الفرح} اللي لمّت، حوالين نفس الترابيزة، عيلتين العروسين.

%%K t04-c1-02-1 p
2. --- إحنا في \VUgras{الخريف}، في \VUgras{شهر سبتمبر}. وهوا \VUgras{أكتوبر}
و\VUgras{نوفمبر} الساقع لسه ما جرّدش الريف من ورقه الأخضر. وهوا الليل دافي
ومعطّر؛ علشان كده العشا بيتعمل تحت \VUgras{الفيرندة} \textsuperscript{(8)} اللي على
الجنينة.

%%K t04-c1-03-1 p
3. --- ومن بعيد، على \VUgras{التلّ} \textsuperscript{(3)}، بيت أهل يوحنا،
\VUgras{قصر} \textsuperscript{(1)} شيك مستخبّي في الخضرة؛ وحقول عنب حلوة، بتطلّع نبيت
ممتاز، مغطّية منحدر التلّ. وبتاع العنب بيزرعها ويعتني بيها.

%%K t04-c1-04-1 p
4. --- وفوق حقول العنب، بيعدّي سرب \VUgras{حجل} \textsuperscript{(2)} اتخضّ من قرب
\VUgras{الخفير} \textsuperscript{(5)}. وده، و\VUgras{بندقيته} تحت باطه و\VUgras{شنطة
الصيد} على كتفه، بيفتّش \VUgras{الشجيرات} \textsuperscript{(4)} مع \VUgras{كلبه}
\textsuperscript{(6)}. بيحرس الأرض ويمنع اللي بيصطادوا خلسة إنهم يفنوا الطرايد.

%%K t04-c1-05-1 p
5. --- وبرّه \VUgras{الفيرندة}، في \VUgras{تعريشة عنب} متسلّقة، متدلّي منها
\VUgras{عناقيد} \textsuperscript{(7)} تخينة.

%%K t04-c1-06-1 p
6. --- وورد الخريف الحلو اللي مزيّن \VUgras{الترابيزة} \textsuperscript{(16)} هو
\VUgras{الأقحوان} \textsuperscript{(9)}؛ و\VUgras{أبو العروسة} \textsuperscript{(10)} حطّ
واحدة
منهم في عروة \VUgras{فراكه}.

%%K t04-c1-07-1 p
7. --- والأكل قرّب يخلص. و\VUgras{الصفرجي} \textsuperscript{(11)} بدأ يجيب أطباق
التحلاية، وحالًا هيشيل حاجات المائدة اللي ما بقاش ليها لزوم:
\VUgras{المعالق} \textsuperscript{(14)}، و\VUgras{الشوك} \textsuperscript{(12)}،
و\VUgras{السكاكين} \textsuperscript{(13)}، و\VUgras{الأطباق الكبيرة} \textsuperscript{(15)}.

%%K t04-tit-3 sub
\VUpk{10.2pt}{40}{القرايب.}
\VUsaut{3.0mm}

%%K t04-c1-08-1 p
8. --- كل الناس مبسوطة، والابتسامة اللي \VUgras{أم} \textsuperscript{(17)} العريس بتبصّ
بيها لولادها بتدلّ على سعادتها.

%%K t04-c1-09-1 p
9. --- الجوازة دي بتوصّل عيلتين غنيّتين من البلد. يوحنا، \VUgras{العريس}
\textsuperscript{(18)}، \VUgras{ابن} صاحب أرض كبير، وهو بقى \VUgras{جوز البنت} بتاعة
صاحب مصنع شاطر.

%%K t04-c1-10-1 p
10. --- \VUgras{العروسين} صغيّرين، ومع كده هما \VUgras{مخطوبين} من زمان.
و\VUgras{العروسة} \textsuperscript{(19)}، مارثا، شكلها حلو أوي في \VUgras{فستانها الأبيض}
المزيّن بزهر البرتقان.

%%K t04-c1-11-1 p
11. --- و\VUgras{حماها} \textsuperscript{(20)} مبسوط أوي إن عنده \VUgras{مرات ابن}
لطيفة كده، و\VUgras{حماة} \textsuperscript{(21)} يوحنا بتبتسم وهي شايفة \VUgras{بنتها}
حلوة أوي كده.

%%K t04-c1-12-1 p
12. --- وفي آخر الترابيزة قاعدين باقي \VUgras{المعازيم} \textsuperscript{(22)}:
\VUgras{القرايب، والأصحاب، وولاد العم، وبنات العم}. و\VUgras{رئيس الصفرجية}
\textsuperscript{(23)}، والفوطة تحت باطه، بيخدمهم بنشاط.

%%K t04-tit-4 sub
\VUpk{10.2pt}{40}{الأصحاب.}
\VUsaut{3.0mm}

%%K t04-c1-13-1 p
13. --- وعلى يمين الترابيزة، أهي \VUgras{آنسة} \textsuperscript{(24)}، \VUgras{صاحبة}
العروسة، وبعدين

%%K t04-c1-13-1 p suite
\VUgras{أخوها}، اللي هو \VUgras{إشبينها} \textsuperscript{(25)}. وهو بيتكلّم مع
\VUgras{إشبينته} \textsuperscript{(26)}، أخت العريس، اللي بالجوازة دي بقت
\VUgras{نسيبة} مارثا. دي بنت صغيرة حلوة. ولابسة \VUgras{مجوهرات} حلوة،
و\VUgras{عقد لؤلؤ}، و\VUgras{أسورة} دهب.

%%K t04-c1-14-1 p
14. --- وجنبهم، الأستاذ اللي واقف ورافع كباية هو \VUgras{صاحب}
\textsuperscript{(27)} للعيلة. وهو واحد من \VUgras{شهود العريس}. وهو \VUgras{بيشرب
نخب} العروسين ويتمنّى لهم سعادة طويلة. لابس بدلة سهرة، فراك
و\VUgras{جزمة لمّاعة} \textsuperscript{(28)}. وهو اللي مرافق \VUgras{تيتة}
\textsuperscript{(29)} اللي \VUgras{أرملة}.

%%K t04-c1-15-1 p
15. --- والشاب اللي لابس \VUgras{فراك} \textsuperscript{(31)} وشنبه الأسود الرفيّع، هو
\VUgras{نسيب} \textsuperscript{(30)} يوحنا. وهو مهندس كبير. وهو قاعد جنب \VUgras{ستّ
شابة} \textsuperscript{(33)} شيك أوي، \VUgras{أخت مارثا الكبيرة} وأم البنوتة الحلوة اللي
جاية، ماسكة إيد ابن عمها، تقدّم وردة و\VUgras{تقول لها مساء الخير}. الستّ دي
لابسة \VUgras{حلق} حلو أوي و\VUgras{طرحة راس} شيك.

%%K t04-c1-16-1 p
16. --- وابن العم الصغير، اللي شكله غريب أوي بـ\VUgras{بلوزته}
\textsuperscript{(36)} و\VUgras{بنطلونه القصير المنفوخ} \textsuperscript{(37)}، حالته تحزن
أوي. هو \VUgras{يتيم} \textsuperscript{(35)}. وهو صغير خالص فقد أبوه وأمه. وعمه بقى
\VUgras{الوصي} \textsuperscript{(38)} عليه، وفضل من غير جواز علشان يربّي الطفل المسكين.
وده راجل طيب. وهو أصلع شوية وقصير النظر أوي؛ وبيستعمل \VUgras{نضّارة}.

%%K t04-tit-5 sub
\VUsaut{2.4mm}
\VUpk{10.2pt}{40}{التحلاية.}
\VUsaut{3.0mm}

%%K t04-c1-17-1 p
17. --- ووراهم، على ترابيزة مغطّاة بـ\VUgras{مفرش}، \VUgras{التحلاية}
\textsuperscript{(39)} محضّرة.

%%K t04-c1-17-1 p suite
في الصينية الكبيرة أهي الفاكهة، \VUgras{خوخ} \textsuperscript{(40)}،
و\VUgras{كمّترى} \textsuperscript{(41)}، و\VUgras{تفاح} \textsuperscript{(42)}، و\VUgras{مشمش}
\textsuperscript{(43)}، و\VUgras{عنب} \textsuperscript{(44)}. وجنبها،
\VUgras{أناناس} \textsuperscript{(45)}، و\VUgras{تورتة} \textsuperscript{(46)} حلوة،
و\VUgras{أزايز مشروبات} \textsuperscript{(48)}، و\VUgras{شمبانيا} \textsuperscript{(49)}،
وكمان أعمدة من \VUgras{الأطباق} \textsuperscript{(47)} النضيفة.

%%K t04-c1-18-1 p
18. --- و\VUgras{الساقي} \textsuperscript{(50)} بيفتح \VUgras{أزازة} \textsuperscript{(52)}
نبيت قديم بـ\VUgras{البريمة} \textsuperscript{(51)}. و\VUgras{الفلّينة}
\textsuperscript{(53)} شكلها صعبة الطلوع أوي. والساقي ده حاطط \VUgras{فوطة}
\textsuperscript{(54)} تحت باطه.

%%K t04-c1-19-1 p
19. --- وبعد العشا، الرجّالة هيروحوا حجرة التدخين، والستّات هيروحوا يحضّروا
نفسهم للرقص.

%%K t04-tit-6 sub
\VUsaut{5.0mm}
\VUcentre{11.4pt}{0}{\textit{المشهد التاني.}}
\VUsaut{1.6mm}
\VUpk{11.4pt}{40}{عيد ميلاد جدو.}
\VUsaut{2.6mm}

%%K t04-tit-7 sub
\VUpk{10.2pt}{40}{الزيارة.}
\VUsaut{3.0mm}

%%K t04-c2-01-1 p
1. --- عدّت عشر سنين، وأهل يوحنا بقوا كبار وبيحبّوا أوي أحفادهم التلاتة،
\VUgras{حفيدين} \textsuperscript{(2)} و\VUgras{حفيدة} \textsuperscript{(3)}. وعلشان النهارده
\VUgras{عيد ميلاد جدو}، الولاد جايين يهنّوهم.

%%K t04-c2-02-1 p
2. --- وبطرس الصغير لسه صغيّر خالص، عنده تلات سنين بس. وهو شايف
\VUgras{القطة} \textsuperscript{(1)} بتقرّب. وعايز يحسّس على \VUgras{فروها} الحرير وعلى
\VUgras{ديلها} الطويل، وهو مش واخد باله إن تحت \VUgras{رجليها} الحلوة مستخبّية
ضوافر شريرة مدبّبة يمكن تخرمشه.

%%K t04-c2-03-1 p
3. --- وأخته جابرييلا، اللي ماسكة إيده، أعقل منه بكتير؛ ومارسيلوس،
بـ\VUgras{بالطوه} \textsuperscript{(4)} الكبير و\VUgras{حزامه} \textsuperscript{(5)} الجلد،
شكله
كبر شوية، بالرغم من بنطلونه القصير اللي مبيّن \VUgras{شرابه}
\textsuperscript{(6)}.

%%K t04-c2-04-1 p
4. --- وأبوهم لابس \VUgras{بالطو} \textsuperscript{(7)} الشتا التخين
بـ\VUgras{ياقة الفرو} \textsuperscript{(8)}، وفي \VUgras{جيبه} \textsuperscript{(9)} عنده
طرحة.
ومراته لابسة برنيطة اللبد بتاعتها المزيّنة بـ\VUgras{الريش} \textsuperscript{(10)}،
و\VUgras{جاكيتتها} \textsuperscript{(11)}، و\VUgras{البوا الفرو} \textsuperscript{(12)}
بتاعها،
و\VUgras{المنشون} \textsuperscript{(13)} بتاعها.

%%K t04-c2-05-1 p
5. --- والدنيا ساقعة أوي، علشان الشتا حكمت. \VUgras{التلج} \textsuperscript{(19)} نزل
طول النهار وسقوف البيوت كلها بيضا. ومع \VUgras{الليل}، البرد بقى أقسى كمان. وفي
السما، \VUgras{القمر} \textsuperscript{(17)} و\VUgras{النجوم} \textsuperscript{(18)} بتلمع
وبتخترق
\VUgras{الضلمة}.

%%K t04-tit-8 sub
\VUsaut{4.2mm}
\VUpk{10.2pt}{40}{المرض.}
\VUsaut{3.0mm}

%%K t04-c2-06-1 p
6. --- و\VUgras{الأعمى} \textsuperscript{(20)} المسكين اللي بيعدّي الميدان قدام
\VUgras{الأجزخانة} \textsuperscript{(16)}، و\VUgras{كلبه البودل} \textsuperscript{(21)}
بيقوده،
حالته تحزن أوي. ويوستينيا، \VUgras{الشغّالة} \textsuperscript{(14)}، بتجيب من المطبخ
\VUgras{كوباية} \textsuperscript{(15)} \VUgras{مغلي أعشاب} سخن أوي ومحلّي كويس.

%%K t04-c2-07-1 p
7. --- و\VUgras{تيتة} \textsuperscript{(23)}، اللي عايزة تدوّر على \VUgras{نضّارتها}
\textsuperscript{(26)} على الترابيزة علشان تقرا \VUgras{الروشتّة} \textsuperscript{(27)} اللي
\VUgras{الدكتور} \textsuperscript{(25)} كتبها حالًا، بتقوم من على الفوتيه بتاعها وهي
متّكية على \VUgras{عصايتها} \textsuperscript{(24)}.

%%K t04-c2-08-1 p
8. --- وهي بتتعب كتير من \VUgras{الوعكات}، و\VUgras{الدوخة}،
و\VUgras{الزكام}، و\VUgras{وجع السنان}؛ ورجليها ضعيفة وعينيها ما بقتش كويسة أوي.
ومع كده، بتحب تهزّر مع \VUgras{دكتورها} القديم اللي عايز على طول يكتب لها
\VUgras{شربة} \textsuperscript{(28)}. والنهارده هي مبسوطة أوي علشان عيلتها الصغيرة
حواليها.

%%K t04-c2-09-1 p
9. --- والحجرة المقفولة كويس مريحة. وفي \VUgras{المدفأة} \textsuperscript{(29)} الرخام
ولعانة

%%K t04-c2-09-1 p suite
\VUgras{نار حلوة أوي} \textsuperscript{(30)}، والواحد بيحسّ بالسعادة في \VUgras{نور}
\VUgras{اللمبة} \textsuperscript{(31)} الهادي اللي ولّعوها حالًا.

%%K t04-tit-9 sub
\VUsaut{4.2mm}
\VUpk{10.2pt}{40}{جدو.}
\VUsaut{3.0mm}

%%K t04-c2-10-1 p
10. --- وحتى \VUgras{جدو} \textsuperscript{(33)} نسي إنه كان عيّان، وإن
\VUgras{الروماتيزم} مثبّته في \VUgras{الفوتيه} \textsuperscript{(34)} بتاعه، وإنه إمبارح
لسه كان عنده \VUgras{سخونية وإغماء}. ما بقاش فاكر إنه الشتا اللي فات تعب من
\VUgras{الالتهاب الرئوي}، وإن \VUgras{كحّة} ناشفة ومتعبة خلّته يتألّم أوي.

%%K t04-c2-10-2 p
ومع كده اتشفى بصعوبة أوي. ومن ساعتها بقى أحسن شوية. بس رجله اللي حاططها على
\VUgras{مخدّة} \textsuperscript{(38)} فوق \VUgras{كرسي صغير} \textsuperscript{(39)}، لسه
بتوجعه هي كمان.

%%K t04-c2-11-1 p
11. --- ولما يبقى ملفوف كويس في \VUgras{روبه} \textsuperscript{(35)} الدافي
بـ\VUgras{زراره} \textsuperscript{(36)} التخينة اللي من صدف، ولما تبقى جنبه، علشان تقرا
له الجرنال، \VUgras{اليتيمة} \textsuperscript{(40)} الصغيرة لابسة \VUgras{طرحة} بيضا
و\VUgras{فستان} \textsuperscript{(42)} أزرق و\VUgras{كاب} \textsuperscript{(41)}، ساعتها ما
بيشتكيش.

%%K t04-c2-12-1 p
12. --- وكل سنة، لما \VUgras{النتيجة} \textsuperscript{(43)} تدلّ تاني على
\VUgras{تاريخ} أول \VUgras{ديسمبر}، بيستنّى \VUgras{أحفاده} بفارغ الصبر، علشان
هو عارف إنهم مش هينسوا \VUgras{التاريخ} المهم ده.

%%K t04-c2-13-1 p
13. --- وبيفتكر بتأثّر الزمن البعيد أوي لما راح بنفسه يهنّي \VUgras{المشير}
الإمبراطوري المجيد، اللي \VUgras{صورته} \textsuperscript{(44)} معلّقة على الحيطة، واللي
هو \VUgras{جدّ جدّ} ولاده.

\VUsaut{6.0mm}
\VUfilet{22mm}
